% ============================================================================
%  SECTION 7 — ETAT DE L'ART : SITUER ISTA ET FISTA       (Membre 4 — Romain)
%
%  NOTE D'INTEGRATION (Membre 1) : ce contenu reprend l'etat de l'art redige par
%  Romain. Il a ete renumerote (section 7, anciennement "0") et aligne sur les
%  macros communes (\prox, \softthr, \norm, \T, \Fopt...). Le fond et la
%  structure d'origine (4 familles + synthese + discussion) sont conserves.
% ============================================================================
\section{Etat de l'art : situer ISTA et FISTA}
\label{sec:etatart}

ISTA et FISTA ne sont qu'un point dans un paysage plus large de methodes
capables de resoudre le probleme composite~\eqref{eq:lasso}. Le critere
d'evaluation du sujet B demande explicitement de comparer les alternatives :
c'est l'objet de cette section. On retient quatre familles --- les methodes de
gradient proximal, la descente par coordonnees, l'ADMM et le Newton proximal ---
et on confronte chacune a ISTA/FISTA selon trois axes juges decisifs en
pratique : la vitesse de convergence (en nombre d'iterations), le cout d'une
iteration, et la structure de probleme pour laquelle la methode est la mieux
adaptee.

\subsection{Le cadre commun : minimiser une somme $f + g$}

Toutes ces methodes attaquent le meme objet : une somme d'une partie lisse $f$
(de gradient $L$-lipschitzien, $L = \norm{A\T A}_2$) et d'une partie convexe non
lisse $g$. Ce qui les distingue, c'est la maniere dont elles exploitent cette
structure :
\begin{itemize}
  \item \textbf{ISTA/FISTA} separent $f$ et $g$ : un pas de gradient sur $f$,
        puis l'operateur proximal de $g$ (le seuillage doux). C'est l'approche
        \emph{full-gradient} : chaque iteration met a jour tout le vecteur $x$.
  \item La \textbf{descente par coordonnees} exploite la separabilite de
        $g = \lambda\norm{\cdot}_1 = \sum_j \lambda |x_j|$ pour ne mettre a jour
        qu'une coordonnee a la fois.
  \item L'\textbf{ADMM} dedouble les variables (une copie pour $f$, une pour $g$)
        et les recoud via un multiplicateur.
  \item Le \textbf{Newton proximal} exploite en plus la courbure de $f$ (la
        hessienne $A\T A$), la ou ISTA/FISTA n'utilisent que le gradient.
\end{itemize}
On reconnait ISTA/FISTA comme le membre le plus \guillemotleft{}~elementaire~\guillemotright{} de cette
famille : ordre 1, separation minimale, aucune structure supplementaire
exploitee. Cette sobriete est a la fois leur force (generalite, simplicite, cout
par iteration faible) et leur limite (convergence sous-lineaire).

\subsection{Descente par coordonnees}

\paragraph{Principe.} Plutot que de bouger toutes les composantes de $x$
ensemble, on les met a jour une par une : on fixe toutes les coordonnees sauf la
$j$-ieme et on minimise $F$ exactement selon $x_j$. Pour le LASSO, ce
sous-probleme a une variable admet une solution explicite --- un seuillage doux
sur le residu partiel :
\[
  x_j \leftarrow \softthr\!\left(
    \frac{1}{\norm{a_j}_2^{2}}\, a_j\T\big(y - A x + a_j x_j\big),\;
    \frac{\lambda}{\norm{a_j}_2^{2}}
  \right),
\]
ou $a_j$ est la $j$-ieme colonne de $A$. On balaie les coordonnees cycliquement
(ou aleatoirement) jusqu'a stabilisation.

\paragraph{Vitesse.} Pour le LASSO, la descente par coordonnees converge a un
taux asymptotiquement \emph{lineaire}, $F(x_k) - \Fopt = O(\rho^{k})$ avec
$\rho < 1$, une fois identifie le support de la solution --- bien mieux que le
$O(1/k)$ d'ISTA et le $O(1/k^2)$ de FISTA dans ce regime final. C'est d'ailleurs
la methode au c\oe ur de \texttt{glmnet}, l'implementation de reference du LASSO
en statistique.

\paragraph{Cout par iteration.} Une mise a jour de coordonnee est tres bon
marche ($O(n)$ si l'on maintient le residu $r = y - Ax$ de maniere
incrementale). Un balayage complet des $p$ coordonnees coute $O(np)$, comparable
a une iteration d'ISTA. L'avantage decisif est l'identification automatique du
support : les coordonnees nulles le restent.

\paragraph{Quand la preferer.} Lorsque la penalite est separable (comme
$\norm{\cdot}_1$) et que les colonnes de $A$ sont peu correlees ; choix par
defaut en grande dimension statistique ($p \gg n$) avec solution tres
parcimonieuse. En revanche, la methode est intrinsequement sequentielle (donc
difficile a paralleliser) et perd son avantage quand $g$ n'est pas separable.

\subsection{ADMM (Alternating Direction Method of Multipliers)}

\paragraph{Principe.} On reformule le probleme en dedoublant la variable via une
copie $z$ et la contrainte $x = z$ :
\[
  \min_{x, z}\; f(x) + g(z) \quad\text{s.c.}\quad x = z,
\]
puis on minimise le lagrangien augmente
$\mathcal{L}_\rho(x,z,u) = f(x) + g(z) + \tfrac{\rho}{2}\norm{x - z + u}_2^2$ en
alternant trois pas :
\[
\begin{aligned}
  x_{k+1} &\leftarrow \argmin_x\; f(x) + \tfrac{\rho}{2}\norm{x - z_k + u_k}_2^2
            && \text{(systeme lineaire),}\\
  z_{k+1} &\leftarrow \prox_{g/\rho}(x_{k+1} + u_k)
            && \text{(seuillage doux),}\\
  u_{k+1} &\leftarrow u_k + x_{k+1} - z_{k+1}
            && \text{(montee de gradient duale).}
\end{aligned}
\]
Le pas en $z$ est exactement le seuillage doux deja implemente pour ISTA.

\paragraph{Vitesse.} En pratique, l'ADMM atteint une precision moderee en
nettement moins d'iterations qu'ISTA (souvent quelques dizaines). Theoriquement,
sa vitesse est $O(1/k)$ pour des objectifs convexes generaux, et lineaire sous
hypotheses fortes ; elle depend toutefois sensiblement du parametre de penalite
$\rho$, dont le reglage est delicat.

\paragraph{Cout par iteration.} C'est le point faible : le pas en $x$ exige de
resoudre un systeme lineaire $(A\T A + \rho I)\,x = b$. Si $A$ est fixe, on
factorise $A\T A + \rho I$ une seule fois (Cholesky, $O(p^3)$) et chaque
iteration ne coute plus que deux substitutions $O(p^2)$ : l'ADMM devient alors
tres competitif. Mais si $\rho$ change ou si $A$ est trop grande pour etre
factorisee, ce pas redevient couteux. Autre nuance : contrairement a ISTA/FISTA
et a la descente par coordonnees, les iteres d'ADMM ne sont pas exactement
parcimonieux avant convergence (seul le bloc $z$ l'est).

\paragraph{Quand le preferer.} Quand le probleme se decompose naturellement
(regularisations multiples, contraintes, problemes distribues), ou quand la
factorisation peut etre amortie sur de nombreuses iterations. L'ADMM brille par
sa modularite : on peut remplacer $g$ par n'importe quelle fonction dont on sait
calculer le prox.

\subsection{Newton proximal}

\paragraph{Principe.} ISTA majore $f$ par un modele quadratique isotrope
$\tfrac{1}{2t}\norm{x - x_k}_2^2$ (une \guillemotleft{}~boule~\guillemotright{}). Le Newton proximal
remplace ce modele par la vraie courbure : il utilise la hessienne
$H = \nabla^2 f = A\T A$ et resout, a chaque pas, un sous-probleme de la forme
\[
  x_{k+1} = \argmin_x\;
    \nabla f(x_k)\T (x - x_k)
    + \tfrac{1}{2}(x - x_k)\T H (x - x_k)
    + g(x),
\]
soit un LASSO local, resolu approximativement par un solveur interne (souvent de
la descente par coordonnees). On peut aussi le voir comme une methode de Newton
semi-lisse appliquee a la condition d'optimalite $0 \in \partial F(\xopt)$.

\paragraph{Vitesse.} C'est le regime le plus rapide pres de l'optimum :
convergence superlineaire, voire quadratique localement
($\norm{x_{k+1} - \xopt} = O(\norm{x_k - \xopt}^2)$). La ou FISTA demande des
centaines d'iterations pour gagner les derniers chiffres de precision, le Newton
proximal les obtient en une poignee de pas.

\paragraph{Cout par iteration.} Tres eleve : chaque iteration exige une
information de second ordre (la hessienne ou son action) et la resolution d'un
sous-probleme interne. Le pari n'est gagnant que si le faible nombre
d'iterations externes compense leur cout unitaire --- ce qui suppose en general
une phase d'amorcage par une methode d'ordre 1 (typiquement FISTA) avant de
basculer sur Newton.

\paragraph{Quand le preferer.} Quand on vise une haute precision sur un probleme
de taille moderee, ou quand $f$ est bien conditionnee et la hessienne
exploitable. A l'inverse, en tres grande dimension ou former $A\T A$ est
prohibitif, les methodes d'ordre 1 restent incontournables.

\subsection{Synthese comparative}

Le tableau~\ref{tab:comparatif} recapitule les compromis. Aucune methode ne
domine partout : le bon choix depend de la precision visee, de la taille du
probleme et de la structure exploitable.

\begin{table}[h]
\centering
\small
\begin{tabular}{@{}lllll@{}}
\toprule
\textbf{Methode} & \textbf{Ordre} & \textbf{Vitesse} &
\textbf{Cout / iteration} & \textbf{Terrain de predilection} \\
\midrule
ISTA            & 1 & $O(1/k)$            & faible ($Ax$, $A\T r$)        & socle simple, general \\
FISTA           & 1 & $O(1/k^2)$          & faible (idem ISTA)            & precision moyenne, par defaut \\
Coord. descent  & 1 & lineaire (asympt.)  & tres faible (1 coord.)        & $g$ separable, $p \gg n$ \\
ADMM            & 1 & $O(1/k)$, $\rho$-dep.& moyen (systeme lin.)         & decomposable, distribue \\
Prox. Newton    & 2 & superlin. / quadr.  & eleve (hessienne + interne)   & haute precision, taille moderee \\
\bottomrule
\end{tabular}
\caption{Comparaison des principales methodes pour le LASSO. \guillemotleft{}~Ordre~\guillemotright{}
designe l'ordre de l'information derivee utilisee. Le regime lineaire de la
descente par coordonnees et le regime quadratique de Newton sont \emph{locaux}
(atteints pres de l'optimum).}
\label{tab:comparatif}
\end{table}

\paragraph{Discussion des compromis.} Trois arbitrages structurent ce paysage.
\begin{itemize}
  \item \emph{Vitesse vs. cout par iteration.} Plus une methode converge en peu
        d'iterations (Newton, ADMM), plus ces iterations sont cheres. ISTA/FISTA
        font le pari inverse : des iterations tres bon marche, quitte a en faire
        beaucoup. En grande dimension, ce pari est souvent le bon.
  \item \emph{Generalite vs. exploitation de structure.} ISTA/FISTA ne demandent
        que $\nabla f$ et le prox de $g$ --- d'ou leur generalite. La descente
        par coordonnees exige la separabilite de $g$ ; le Newton proximal une
        courbure exploitable ; l'ADMM une decomposition commode. Ces hypotheses,
        quand elles tiennent, paient.
  \item \emph{Memoire et robustesse.} ISTA/FISTA et la descente par coordonnees
        ont une empreinte memoire minimale. L'ADMM stocke une variable duale et
        idealement une factorisation ; le Newton proximal manipule de
        l'information de second ordre, le plus gourmand. Cote robustesse, ISTA
        est le plus previsible ; FISTA introduit des oscillations non monotones
        (cf.~section~\ref{sec:exp}) que des variantes \emph{restart} corrigent ;
        l'ADMM est sensible a $\rho$ ; Newton depend d'un bon amorcage.
\end{itemize}

\paragraph{Position d'ISTA/FISTA.} En definitive, FISTA occupe un point
d'equilibre remarquable : sans hypothese de structure au-dela de la separation
$f + g$, sans parametre delicat a regler, avec un cout par iteration minimal, il
atteint le meilleur taux possible pour une methode du premier ordre
($O(1/k^2)$, optimal au sens de Nesterov). C'est ce qui en fait le \guillemotleft{}~cheval de
bataille~\guillemotright{} par defaut, et la brique de base a partir de laquelle les autres
methodes se construisent ou se comparent --- d'ou sa place centrale dans ce
projet, les autres approches servant de points de repere pour en mesurer les
forces et les limites.
